\section*{Resumen}\addcontentsline{toc}{section}{Resumen}
Se implementaron seis metaheurísticas para buscar ciclos hamiltonianos en instancias del FHCP Challenge Set: algoritmo genético (GA), recocido simulado (SA), Ant System (AS, de la familia ACO), colonia artificial de abejas (ABC), evolución diferencial (DE) y optimización por enjambre de partículas (PSO). El problema se representó mediante permutaciones de vértices y un costo que cuenta las aristas ausentes, incluido el cierre de la ruta. Todo éxito exigió validación independiente y evidencia de haber sido encontrado dentro del presupuesto.

El trabajo conserva protocolos distintos. La campaña inicial de benchmarking completó 1\,044 procesos, con cero éxitos en sus 900 intentos principales de evaluación y escala. Después se incorporaron inicialización voraz, propuestas guiadas por aristas y optimizaciones específicas. La evaluación distribuida registró \TotalAttempts{} intentos en \StartedGraphs{} grafos: \TotalSuccesses{} éxitos por intento y \SolvedGraphs{} grafos distintos certificados, con \TotalErrors{} incidencias conservadas. Se conciliaron 60 combinaciones en \CompleteGraphs{} grafos y cinco quedaron parciales.

En los grafos con cobertura completa, ABC produjo 116 éxitos en 1\,440 intentos y alcanzó los 15 grafos certificados, la mayor cobertura observada entre las familias. Esta conclusión es descriptiva y está limitada a la selección, parámetros y presupuesto evaluados. Los tiempos contienen censura por el plazo y proceden de equipos distintos; no permiten atribuir una superioridad general ni comparar aceleraciones con publicaciones externas.

Las exploraciones complementarias quedaron cerradas con dos grafos adicionales en la variante guiada y cinco en la variante con reparación v4. Esta última registró \RepairAttempts{} intentos, \RepairSuccesses{} éxitos y \RepairErrors{} interrupciones, con 30 segundos y cuatro procesos. La unión de las tres entregas contiene \UnionCertifiedGraphs{} grafos distintos certificados: supera la referencia aproximada de 20, sin alcanzar la meta original de 60. Las tasas y tiempos de protocolos diferentes se conservan separados; las combinaciones no ejecutadas quedan fuera del cierre. El informe presenta tablas y figuras reproducibles, comparación con fuentes primarias y límites de interpretación.

\textbf{Palabras clave:} ciclo hamiltoniano; metaheurísticas; FHCP; búsqueda guiada; reproducibilidad; evaluación experimental.
\clearpage
